%%% Jim Ashe (jim.ashe@wgu.edu; x4209)%%%%
\documentclass{exam}
\usepackage{WGU_worksheet_preamble}

%%%%%%%%%%%%%%%%%%%%%%%%%%%%%%%%%%%%%%%%%%
%%header%%%%%%%%%%%%%%%%%%%%%%%%%%%%%%%%%%
\headrule
\header{ 
    \href{mailto:matt.alexander@wgu.edu?subject=C960: probability worksheet}{Dr.\,Matt Alexander}, 
    \href{mailto:joshua.bilbrey@wgu.edu?subject=C960: probability worksheet}{Dr.\,Josh Bilbrey}, and
    \href{mailto:malcolm.rupert@wgu.edu?subject=C960: probability worksheet}{Dr.\,Malcom Rupert}
    \\ \textbf{C960: Discrete Probability}}
{Practice Problems}
{}

%%%%%%%%%%%%%%%%%%%%%%%%%%%%%%%%%%%%%%%%%
\lfoot{\footnotesize Suggestions or errors?\\ Contact the author:   \href{mailto:jim.ashe@wgu.edu?subject=C960: probability worksheet}{Dr.\,Jim Ashe}.}
\cfoot[]{Page \thepage\ of \numpages}
\rfoot{\today}
%%%%%%%%%%%%%%%%%%%%%%%%%%%%%%%%%%%%%%%%%
\newcommand{\itemSpace}{36pt} %adjusts spacing between questions. See question environment
\usepackage{color}
%\printanswers
\shadedsolutions
\colorlet{SolutionColor}{lightBlueWGU!15!white}

%\definecolor{SolutionColor}{gray}{0.8}
\begin{document}

%%%%%Advanced Counting%%%%%%%%%%%%%%%%%%%
\section*{Advanced Counting} 
\begin{questions}
\setlength\itemsep{\itemSpace} %set length above question. Needs to be set for each question environment
\question A college has $1105$ students. This semester $320$ students are
taking a math course, $405$ of the students are taking a history
class, and $250$ of the students taking a math class are not enrolled
in a history class.
    \begin{parts}
        \part Construct a Venn diagram, and write the appropriate numbers in the four regions
            \begin{solution}
            
            \begin{tikzpicture}[scale=1.0]
            	\def\circleA{(0,0) circle (1.5)} %A
            	\def\circleB{(0:2) circle (1.5)} %B
            
            	\colorlet{circle edgeA}{black!75} 
            	\colorlet{circle areaA}{red!100}
            	\colorlet{circle edgeB}{black!75} 
            	\colorlet{circle areaB}{green!100}
            
            	\tikzset{
            		filledA/.style={fill=circle areaA, draw=circle edgeA, thick},
            		filledB/.style={fill=circle areaB, draw=circle edgeB, thick},
            		outlineA/.style={draw=circle edgeA, thick},
            		outlineB/.style={draw=circle edgeB, thick},
            	}
            
            	\begin{scope}[fill opacity=0.5]         
            			%\fill[filledB] \circleB;    
            			%\fill[filledA] \circleA;
            	\end{scope}
            
            	\draw[outlineA] \circleA node {$M$} node[below=10pt] {$250$}; 
            	\draw[outlineB] \circleB node {$H$} node[below=10pt] {$335$}; 
            	\node at (0:1) {$70$}; %%intersection node%%
            	\node[anchor=south] at (current bounding box.north) {$450$}; 
            	\draw[thick] (-2,-2) rectangle (4,2) node[anchor=south east] at (current bounding box.south east) {$U$}; 
            \end{tikzpicture} 
            
            $n(M\cap H)=320-250$=70, $n(H-M)=405-70=335$, and $n\left((M\cup H)'\right)=1105-250-70-335=450$.
            \end{solution}
        \part How many of the students are taking math and history?
            \begin{solution}
            $n(H\cap H)=70$
            \end{solution}
        \part How many students are taking history but not math?
            \begin{solution}
            $n(H-M)=335$
            \end{solution}
        \part How many students are taking math or history?
            \begin{solution}
            $n(M\cup H)=250+70+335=655$
            \end{solution}
        \part How many students are taking neither math nor history?
            \begin{solution}
            $n\left((M\cup H)'\right)=n(U)-n(M\cup H)=1105-655=450$
            \end{solution}
    \end{parts}


\question \emph{Survivor}, \emph{CSI}, and \emph{ER} all belong to the Thursday
night prime time lineup. A group of $600$ college students were surveyed.
$350$ students watch \emph{Survivor}, $280$ watch \emph{CSI}, and
$320$ watch \emph{ER}. $50$ students watch all three shows. $190$
students watch both \emph{Survivor} and \emph{ER}. $180$ students
watch both \emph{Survivor} and \emph{CSI}. $50$ students watch only
\emph{ER}.
    \begin{parts}        
    \part Construct a Venn diagram, and write the appropriate numbers in the eight regions.
        \begin{solution}
            
            \begin{tikzpicture}[scale=.6]    
            \def\firstcircle{(90:1.75cm) circle (2.5cm)} 
            \def\secondcircle{(210:1.75cm) circle (2.5cm)} 
            \def\thirdcircle{(330:1.75cm) circle (2.5cm)}
            \begin{scope}     
            	\fill[white]\firstcircle;       
            	\fill[white] \secondcircle;       
            	\fill[white] \thirdcircle;     
            \end{scope}     
            \begin{scope}
            	\clip 
            	\secondcircle;       
            	\fill[white] 
            	\firstcircle;     
            \end{scope}     
            \begin{scope}       
            	\clip 
            	\secondcircle;       
            	\fill[white] 
            	\thirdcircle;     
            \end{scope}     
            \begin{scope}       
            	\clip 
            	\firstcircle;       
            	\fill[white] 
            	\thirdcircle;     
            \end{scope}     
            \begin{scope}       
            	\clip \firstcircle;       
            	\clip \secondcircle;       
            	\fill[white] \thirdcircle;     
            \end{scope}     
            \draw 
            	\firstcircle node[text=black] {$S$} node[above=10pt] {$30$};     
            \draw 
            	\secondcircle node [text=black,below left] {$CSI$} node[below left=10pt] {$20$};     
            \draw 
            	\thirdcircle node [text=black,below right] {$ER$} node[below right=10pt] {$50$};     
            \begin{scope}[font=\small]       
            	\node(0,0) {$50$}; %AnBnC       
            	\draw(30:1.5cm) node {$140$}; %AnC       
            	\draw(150:1.5cm) node {$130$}; %AnB       
            	\draw(270:1.5cm) node {$80$}; %BnC    
            \end{scope}
            
            \draw[thick] (-4.8,-3.55) rectangle (4.75,4.4) node[anchor=south east] at (current bounding box.south east) {$U$} node[anchor=north west] at (current bounding box.north west) {$100$}; 
            
            \end{tikzpicture}
        
        Number in $S$ and $ER$ but not $CSI$ is: $190-50=140$, Number
        in $S$ and $CSI$ but not $ER$ is: $180-50=130$. Number in $CSI$
        and $ER$ but not $S$ is: $320-50-140-50=80$. Number only in $S$
        is: $350-130-50-140=30$. Number only in $CSI$ is: $280-130-50-80=20$.
        \end{solution}
    \part How many students watch only \emph{Survivor}?
        \begin{solution}
        From diagram:$30$
        \end{solution}
    \part How many students watch \emph{CSI }or \emph{Survivor}?
        \begin{solution}
        From diagram: $30+130+140+50+80+20=450$, or by the union rule:
        $350+280-180=450$.
        \end{solution}
    \part How many students watch \emph{CSI }or \emph{ER}?
        \begin{solution}
        From diagram: $20+130+50+80+140+50=570$, or by the union rule:
        $280+320-130=470$
        \end{solution}
    \part How many students watch only \emph{CSI}?
        \begin{solution}
        From diagram: $20$
        \end{solution}
    \end{parts}


\question A certain model of pickup truck is available in five exterior colors,
three interior colors, and three interior styles. In addition, the
transmission can be either manual or automatic, and the truck can
have either two-wheel or four-wheel drive. How many different versions
of the pickup truck can be ordered?
    \begin{solution}
        We have five categories, $\underline{\,\,\,\,\,}\times\underline{\,\,\,\,\,}\times\underline{\,\,\,\,\,}\times\underline{\,\,\,\,\,}\times\underline{\,\,\,\,\,}$.
        Filling in the number of choices for each category we have, 
        \[
        \underline{5}\times\underline{3}\times\underline{3}\times\underline{2}\times\underline{2}=180
        \]
    \end{solution}


\question The $\Delta \Delta \Delta$ fraternity must elect a president, vice president, and treasurer from its 30 members. How many ways can this be done?
    \begin{solution}
        There are no repeats and order does matter (since order will determine
        who is pres., vice pres., and treasurer). So we use permutations,
        \begin{eqnarray*}
        _{30}P_{3}=\frac{30!}{(30-3)!}=\frac{30!}{27!} & = & 30\cdot29\cdot28=24,360
        \end{eqnarray*}
    \end{solution}


\question The local branch of a national business chain has 27 employees. If they all shake each others hand at the company picnic, at least how many handshakes occurred?
    \begin{solution}
        There are no repeats and order does NOT matter (Joe shaking hands
        with Betty is the same as Betty shaking hands with Joe). So we use
        combinations,
        \[
        _{27}C_{2}=\frac{27!}{(27-2)!2!}=\frac{27!}{25!2!}=\frac{27\cdot26}{2}=351
        \]
    \end{solution}


\question Using a standard deck of $52$ cards,
    \begin{parts}
    \part how many ways can five cards be drawn?
        \begin{solution}
            Here order matters, e.g., $\{7\heartsuit,K\spadesuit,2\spadesuit,3\clubsuit,3\diamondsuit\}$
            is a different way to draw the cards then $\{2\spadesuit,3\clubsuit,3\diamondsuit,7\heartsuit,K\spadesuit,\}$
            since the order is different. So we use permutations,
            \[
            _{52}P_{5}=\frac{52!}{(52-5)!}=\frac{52!}{47!}=311,875,200
            \]
        \end{solution}
    \part how many five card hands can be dealt?
        \begin{solution}
            Here order does not matter, e.g., $\{7\heartsuit,K\spadesuit,2\spadesuit,3\clubsuit,3\diamondsuit\}$
            is the same hand as 
            $$\{2\spadesuit,3\clubsuit,3\diamondsuit,7\heartsuit,K\spadesuit\}$$
            So we use combinations, 
            \[
            _{52}C_{5}=\frac{52!}{(52-5)!5!}=2,598,960
            \]
        \end{solution}
    \end{parts}    


\question A three-person committee must be chosen from a group of ten men and twelve women. How many committees are possible if it must consist of the following?
    \begin{parts}
    \part Three men.
        \begin{solution}
            No repeats and order doesn't matter$\Rightarrow$combination.
            $_{10}C_{3}=\frac{10!}{(10-3)!3!}=\frac{10\cdot9\cdot8}{3!f}=\frac{720}{6}=120$
        \end{solution}
    
    \part Three women.
        \begin{solution}
            No repeats and order doesn't matter$\Rightarrow$combination.
            $_{12}C_{3}=\frac{12!}{(12-3)!3!}=\frac{12\cdot11\cdot10}{3!}=\frac{1320}{6}=220$
        \end{solution}
    
    \part One woman and two men.
        \begin{solution}
            Here we have two categories: $\underline{\,\,\,\,\,}\times\underline{\,\,\,\,\,}$.
            For the first categories we have $12$ choices (equivalently, $_{12}C_{1}=\frac{12!}{(12-1)!1!}=12$),
            and for the second we have $_{10}C_{2}=\frac{10!}{(10-2)!2!}=\frac{10\cdot9}{2}=45$.
            \\
            So we have $12\times45=540$ possibilities.
        \end{solution}
    
    \part One man and two women, and one of these is selected as a spokesperson of the committee.
        \begin{solution}
            Similarly to above, we have two categories. For the first categories
            we have $10$ choices, and for the second we have $_{12}C_{2}=\frac{12!}{(12-2)!2!}=\frac{12\cdot11}{2}=66$.
            Then we select one of the three to be a spokesperson. \\
            So we have $10\times66\times3=1980$ possibilities.
            
            Equivalently, use a permutation where the order of $2$ members
            does not matter, $\frac{\left( 10\times\frac{12!}{10!}\times 3 \right) }{2!}=1980$
            (like counting the ways to arrange the letters in ``SCHOOL''; the
            order of the ``O'''s does not mater).
        \end{solution}
    \end{parts}

    
\question Each student at State University has a student ID number consisting
of five digits (the first digit is nonzero, and the digits can be
repeated) followed by two of the letters $A$, $B$, $C$, and $D$
(letters cannot be repeated). How many different student numbers are
possible?
    \begin{solution}
        We have five numbers $0-9$ which can be repeated, but the the
        first is nonzero: $9\times10\times10\times10\times10=9\cdot10^{4}=90,000$
        possibilities, and then two of four letters which cannot be repeated.
        It is implied that the order of the letters does matter since $10000AB$
        and $10000BA$ would be two different IDs. So for the letters we have,
        $_{4}P_{2}=\frac{4!}{(4-2)!}=\frac{24}{2}=12$ possibilities. For
        a total of $9\times10^{4}\times12=1,080,000$ possible ID numbers.
    \end{solution}


\question How many ways are there to place 6 identical objects into 3 different bins? 
\begin{solution}
This problem is similar to \#33 and \#35 on the pre-assessment. In general, the number of ways to place $n$ indistinguishable object into $k$ groups is $${n+k-1 \choose n} = {n+k-1 \choose k-1}$$ This is often called a "star and bar" approach. We can line up our identical objects (stars) and place dividers between them (bars) to indicate which bin they will go into. For example, $$\star \star | \star \star | \star \star$$ places 2 objects in each bin, and $$\star \star  \star | \star | \star \star |$$ places 3 into the first two bins and none in the last. Note that we need to count the ways to place $3-1=2$ bars in $6+2$ spots. Order does not matter: $${6+(3-1) \choose (3-1)} = {8 \choose 2}=28$$ Similarly, we can count the number of ways to place 6 stars in the eight places: $${8 \choose 6} = 28$$
\end{solution}


\question How many ways can we add three non-negative integers such that they sum to 9? (order matters)
\begin{solution}
This problem may be more difficult than what you'll see on the assessment. However, I thought it worth including as "counting by bijections" are listed in the assessment competencies methods. Furthermore, it's really no different than \#33 and \#35 on the pre-assessment\\ 

The problem is asking us to count the number of ordered sets of non-negative integers, $(a,b,c)$, such that $$a+b+c=9$$ We need to describe the situation in terms of indistinguishable objects (as above). So we'll use a "star and bar" approach. Only we'll use ones and zeros instead (using bars for $0$'s would get confusing!). \\

Think of $a,b,c$ as bins and the ones as indistinguishable objects. So we will need $9$ ones and $3-2$ zeros. For example $$11101101111$$ corresponds to $3+2+3=9$ and $$11111001111$$ corresponds to $5+0+4$. So now we have a bijection between the solutions to $a+b+c=9$ and these sequences of ones and zeros. And the problem is equivalent to counting the number of ways to choose 2 out 11 objects $${11 \choose 2}=55$$
\end{solution}


\question Which of the following functions maps a $11-$to$-1$ correspondence from the set ${1,2,3,\dots, 44}$ to the set ${1,2,\dots , 11}$?
    
    \begin{oneparchoices}
    \choice $f(x)={x \choose 3}$
    \choice $f(x)=\lfloor \frac{x}{11} \rfloor$
    \CorrectChoice $f(x)=\lceil \frac{x}{11} \rceil$
    \choice $f(x)=\lbrack \frac{x}{11} \rbrack$
    \end{oneparchoices}
    \begin{solution}
        The notation $\lceil x \rceil$ is called a \textit{ceiling} function. It is the smallest integer less than or equal to $x$. For example, $\lceil \frac{1}{11} \rceil=\lceil .09\overline{09} \rceil=1$. So then,
        $$\lceil \frac{1}{11} \rceil=\lceil \frac{2}{11} \rceil=\lceil \frac{3}{11} \rceil=\cdots \lceil \frac{11}{11} \rceil=1$$
        
        Two options here make no sense, ${x \choose 3}$ and $\lbrack \frac{x}{11} \rbrack$. Our only two real options are $f(x)=\lfloor \frac{x}{11} \rfloor$ and $f(x)=\lceil \frac{x}{11} \rceil$. It is to check that the former is incorrect: $\lfloor \frac{1}{11} \rfloor=0$, but $0$ is not in the set ${1,2,\dots , 11}$.
    \end{solution}


\question Which of the following functions maps a $11-$to$-1$ correspondence from $\{1,2,3,\dots, 44\}$ to $\{0,1,\dots , 10\}$?
    
    \begin{oneparchoices}
    \choice $f(x)={x \choose 3}$
    \CorrectChoice $f(x)=\lfloor \frac{x}{11} \rfloor$
    \choice $f(x)=\lceil \frac{x}{11} \rceil$
    \choice $f(x)=\lbrack \frac{x}{11} \rbrack$
    \end{oneparchoices}
    \begin{solution}
        Again, two options here make no sense: ${x \choose 3}$ and $\lbrack \frac{x}{11} \rbrack$. Our only two real options are $f(x)=\lfloor \frac{x}{11} \rfloor$ and $f(x)=\lceil \frac{x}{11} \rceil$. What is different here is the range set. We need to map $11$ numbers to $0$. It is to check that which function does $\lfloor \frac{1}{11} \rfloor=0$. There are only two options. So carefully check that you are getting the correct one. 
    \end{solution}

\question What is the coefficient of $x^{6}y^{4}$ in $(x-2y)^{10}$?
    \begin{solution}
        Recall the Binomial Theorem:
        $$(a+b)^{n}=\sum^{n}_{k=0} {c \choose k}a^{k}b^{n-k}$$
        $a=x$ and $b=-2y$ (be careful with the sign!) \\
        $$(x-2y)^{10}= \sum^{10}_{k=0} {10 \choose k}x^{k}(-2y)^{10-k}$$ 
        Looking at the provided variables $x^{6}y^{4}$. So $k=6$. Now we just plug this in: 
        $${10 \choose 6}x^{6}(-2y)^{10-6}=210(1)^{6}(-2)^{4}=3360$$ 
    \end{solution}

\question What is the value of 
${10 \choose 0}-{10 \choose 1}+{10 \choose 2}+\cdots+{10 \choose 10}$?
    \begin{solution}
        0. Recall the Binomial Theorem:
        $$(a+b)^{n}=\sum^{n}_{k=0} \binom{n}{k} a^{k}b^{n-k}$$ Now rewriting our situation, 
        $${10 \choose 0}-{10 \choose 1}+{10 \choose 2}+\cdots+{10 \choose 10}=\sum^{10}_{k=0} {10 \choose k}(-1)^{k}(1)^{10-k}$$
        If $a=-1$ and $b=1$ then $(-1+1)^{10}=0$. The text also gives this as an identity in 4.19 $$\binom{n}{0}-\binom{n}{1}+\binom{n}{2}-\cdots (-1)^{n}\binom{n}{n}=0$$  
    \end{solution}
\question Suppose that $B=\{0,1,2,\dots 8 \}$. If $f:A\rightarrow B$ is 3-to-1 mapping of $A$ onto $B$, what is $|A|$?

    \begin{solution}
        Recall that $|A|$ notates the number or elements in the set $A$. $B$ has 9 elements, $|B|=9$. $f$ maps 3 elements from $A$ to every element in $B$. So the $|A|=27$.
    \end{solution}


%%%another ? a closed form question??

\end{questions}
\newpage

%%%%%Discrete Probability%%%%%%%%%%%%%%%%
\section*{Discrete Probability}
\begin{questions}
\setlength\itemsep{\itemSpace} %set length above question. Needs to be set for each question environment
\question A card is dealt from a well-shuffled, standard deck of $52$ cards.
Find the probability of being dealt each of the following.
    \begin{parts}
        \part A red card
            \begin{solution}
                $\frac{n(\mbox{red cards})}{n(\mbox{cards})}=\frac{26}{52}=\frac{1}{2}$
            \end{solution}
        \part A queen
            \begin{solution}
                $\frac{n(\mbox{queens})}{n(\mbox{cards})}=\frac{4}{52}=\frac{1}{13}$
            \end{solution}
        \part A club
            \begin{solution}
                $\frac{n(\mbox{clubs})}{n(\mbox{cards})}=\frac{13}{52}=\frac{1}{4}$
            \end{solution}
        \part The queen of clubs
            \begin{solution}
                $\frac{n(\mbox{queen of clubs})}{n(\mbox{cards})}=\frac{1}{52}$
            \end{solution}
        \part A queen or a club
            \begin{solution}
                $\frac{n(\mbox{queens or clubs})}{n(\mbox{cards})}=\frac{4+13-1}{52}=\frac{16}{52}=\frac{4}{13}$
            \end{solution}
        \part Not a queen
            \begin{solution}
                $1-\frac{4}{52}=\frac{48}{52}=\frac{12}{13}$
            \end{solution}
    \end{parts}


\question Three fair coins are tossed. Find each of the following
    \begin{parts}
        \part The sample space
            \begin{solution}
                \forestset{  
                	L1/.style={rectangle, rounded corners, very     thick},   	
                	L2/.style={rectangle, rounded corners, very     thick,
                		edge={black,line width=.75pt,line cap=round}},   	
                	L3/.style={rectangle, rounded corners, very     thick,
                		edge={black,line width=.75pt,line cap=round}},   	
                	L4/.style={rectangle, rounded corners, very     thick,
                	edge={black,line width=.75pt,line cap=round}}, 
                }
                
                \begin{forest}     
                	for tree={
                    		scale=.75,        
                    		grow=0,
                    		reversed, %tree direction         
                    		parent anchor=east,
                    		child anchor=west, %edge anchors   		edge={line cap=round},
                    		outer sep=+1pt, %edge/node connection  		rounded corners,
                    		%minimum width=15mm,
                    		%minimum height=8mm, %node shape         l 		sep=10mm %level distance     
                    		}   
                    [Start,L1
                    	[H,L2         
                    		[H,L3            
                    			[H \{HHH\},L4] 
                    			[T \{HHT\},L4]         
                    		]         
                    		[T,L3          
                    			[H \{HTH\},L4]
                    			[T \{HTT\},L4]
                    		]               
                    	]
                    	[T,L2         
                    		[H,L3            
                    			[H \{THH\},L4,]
                    			[T \{THT\},L4,]         
                    		]         
                    		[T,L3          
                    			[H \{TTH\},L4]
                    			[T \{TTT\},L4]
                    		]               
                    	]  
                    ] 
                \end{forest}
            \end{solution}
        \part The event, $E_{1}$, that exactly two are tails.
            \begin{solution}
                $E_{1}=\{HTT,THT,TTH\}$
            \end{solution}
        \part The event, $E_{2}$, that at least two are tails.
            \begin{solution}
                $E_{2}=\{HTT,THT,TTH,TTT\}$
            \end{solution}
        \part The probability of $E_{1}$
            \begin{solution}
                $p(E_{1})=\frac{3}{8}$
            \end{solution}
        \part The probability and odds of $E_{2}$
            \begin{solution}
                $p(E_{2})=\frac{4}{8}=\frac{1}{2}$ 
            \end{solution}
    \end{parts}


\question Five cards are dealt (without replacement) from a well-shuffled $52$ card standard
deck. Find the probability of the following:
    \begin{parts}
        \part being dealt $5$ spades.
            \begin{solution}
                $\frac{_{13}C_{5}}{_{52}C_{5}}\approx 0.0004951$
            \end{solution}
        \part being dealt $4$ spades and $1$ heart.
            \begin{solution}
                $\frac{_{13}C_{4}\cdot_{13}C_{1}}{_{52}C_{5}}=\frac{_{13}C_{4}\cdot13}{_{52}C_{5}} \approx 0.003576$
            \end{solution}
        \part being dealt all the same suit.
            \begin{solution}
                $\frac{_{4}C_{1}\cdot_{13}C_{5}}{_{52}C_{5}}=\frac{4\cdot_{13}C_{5}}{_{52}C_{5}} \approx 0.00198$
            \end{solution}
        \part not being dealt all the same suit
            \begin{solution}
                This is just $1$ minus the answer above: $1-\frac{4\cdot_{13}C_{5}}{_{52}C_{5}} \approx 0.99801$
            \end{solution}
    \end{parts}


\question A drawer has $7$ socks. $4$ socks are black and $3$ are white socks. John randomly pulls out $4$ socks. Find the probability of the following:
    \begin{parts} 
        \part all $4$ socks are black.
            \begin{solution}
                $\frac{_{4}C_{4}}{_{7}C_{4}}=\frac{1}{35}$
            \end{solution}
        \part exactly $2$ are white.
            \begin{solution}
                $\frac{_{3}C_{2}\cdot_{4}C_{2}}{_{7}C_{4}}=\frac{3\cdot6}{35}=\frac{18}{35}$ \\
                $2$ are white so $2$ are white.
            \end{solution}
        \part at least $3$ are white.
            \begin{solution}
                $\frac{_{3}C_{3}\cdot_{4}C_{0}+_{3}C_{4}\cdot_{4}C_{0}}{_{7}C_{4}}=\frac{1\cdot1+0\cdot1}{35}=\frac{1}{35}$ \\
                $3$ OR $4$ are white.
            \end{solution}
        \part at most $2$ are black. 
            \begin{solution}
                $\frac{_{4}C_{0}\cdot_{3}C_{4}+_{4}C_{1}\cdot_{3}C_{3}+_{4}C_{2}\cdot_{3}C_{2}}{_{7}C_{4}}=\frac{1\cdot0+4\cdot1+6\cdot3}{35}=\frac{22}{35}$ \\
                $0$ or $1$ or $2$ are black. 
            \end{solution}
            \begin{solution}
                Note: Above, $_{3}C_{4}=0$ as there are $0$ ways to choose $4$
                white socks if there are only $3$. I will \textbf{not} have a situation
                like this on the exam.
            \end{solution}
    \end{parts}


\question In a certain game, a player bets \$5. There is a 25\% chance that the player wins \$2, a 25\% chance the player wins \$7, and a 50\% chance the player loses \$5. What is the expected outcome of this game? Round to the nearest cent. 
    \begin{solution} Recall the definition of \emph{expected value}.
        The \emph{expected value} of a discrete random variable $X$ is $$E[X]=\sum_{s\in S}X(s)p(s)$$ where $S$ is the set of possible outcomes of $X$, $p(s)$ is the probability of outcome $s$, and $X(s)$ is the value of $X$ for $s$.
    \end{solution}
    \begin{solution}
        $E[X]=2(.25)+7(.25)+(-5)(.50)=-.25$
    \end{solution}

\question WGU is holding a raffle to raise money for a new internet service.
The tickets cost $\$10$ each. The value and number of each prize
to be given away is listed in the provided table. Find the expected
value of a ticket under the given circumstances. \\
    \begin{center}
    \begin{tabular}{|c|c|c|}
        \hline 
        \textbf{Prize} & \textbf{Value} & \textbf{\# }\tabularnewline
        \hline 
        Vacation & $\$3,000$ & 1\tabularnewline
        \hline 
        Stereo & $\$500$ & 5\tabularnewline
        \hline 
        Books & $\$100$ & 10\tabularnewline
        \hline 
        T-shirt & $\$20$ & 100\tabularnewline
        \hline 
    \end{tabular}
    \end{center}
    \begin{parts}
            \begin{solution} Recall the definition of \emph{expected value}.
                The \emph{expected value} of a discrete random variable $X$ is $$E[X]=\sum_{s\in S}X(s)p(s)$$ where $S$ is the set of possible outcomes of $X$, $p(s)$ is the probability of outcome $s$, and $X(s)$ is the value of $X$ for $s$.
            \end{solution}
            
        \part $200$ tickets
            \begin{solution}
                There are $200-116=84$ losing tickets.\\
                $E[X]=2990(\frac{1}{200})+490(\frac{5}{200})+90(\frac{10}{200})+10(\frac{100}{200})-10(\frac{84}{200})$\\
                \hspace{0cm}\hphantom{$EV$}$=\$32.5$
            \end{solution}
        \part $1000$ tickets
            \begin{solution}
                There are $1000-116=884$ losing tickets.\\
                $E[X]=2990(\frac{1}{1000})+490(\frac{5}{1000})+90(\frac{10}{1000})+10(\frac{100}{1000})-10(\frac{884}{1000})$\\
                \hspace{0cm}\hphantom{$EV$}$=-\$1.5$
            \end{solution}
    \end{parts}


\question A fair six-sided dice is rolled 5 times. Find the probability of the
following:
    \begin{parts}
        \part All $2$'s are rolled.
            \begin{solution}
                $p(\mbox{five }2's)=\frac{1}{6}\cdot\frac{1}{6}\cdot\frac{1}{6}\cdot\frac{1}{6}\cdot\frac{1}{6}=\frac{1}{6^{5}}$
            \end{solution}
        \part Four $2$'s are rolled and the \nth{5} roll is a $1$.
            \begin{solution}
                $p(\mbox{four }2's\cap\mbox{one }1)=(\frac{1}{6}\cdot\frac{1}{6}\cdot\frac{1}{6}\cdot\frac{1}{6})\cdot\frac{1}{6}=\frac{1}{6^{5}}$
            \end{solution}
        \part Not all $2$'s are rolled.
            \begin{solution}
                $1-\frac{1}{6^{5}}$
            \end{solution}
    \end{parts}

\end{questions}
\newpage

%%%%%%Bayes' Theorem%%%%%%%%%%%%%%%%%%%%%
\section*{Baye's Theorem}
\begin{questions}
\setlength\itemsep{\itemSpace} %set length above question. Needs to be set for each question environment
    \begin{solution}
        Recall \textbf{Bayes' Theorem:} For events $X$ and $F$ from a common sample space where $P(X),P(F)\neq 0$, 
        \begin{equation}\label{eqn: bayes1}
            P(F|X)=\frac{P(X|F)P(F)}{P(X)}\tag{version 1}
        \end{equation}

        and equivalently,
        \begin{equation}
            P(F|X)=\displaystyle \frac{P(X|F)P(F)}{P(X|F)P(F)+P(X|\overline{F})P(\overline{F})}\tag{version 2}
        \end{equation}
        While the first version does not appear in ant textbook or pre-assessment practice problems, this equivalent definition is listed in the lesson summary. 
        
        Note that $P(F\cap X)=P(X|F)P(F)$. So version 1 is often shown as:
        \begin{equation}\label{eqn: bayes1}
            P(F|X)=\frac{P(X\cap F)}{P(X)}\tag{version 3}
        \end{equation}
    \end{solution}
    
\question Bag \#1 contains 10 blue marbles. Bag \#2 appears identical but contains 5 red marbles and 5 blue marbles. Suppose you randomly choose a bag and then randomly pick a blue marble. What is the probability that you choose bag \#2?
    \begin{solution}
        We know that a blue marble was chosen. So we are looking for the probability of selecting Bag \#2 \emph{given} that a blue marble was picked, e.g., $P(\#2|B)$. The conditional statement indicates we should use Bayes' Theorem. We'll use first version (seen above):
            $$P(\#2|B)=\frac{P(B|\#2)P(\#2)}{P(B)}$$
        So we need to find the following:
        \begin{itemize}
            \item The probability of selecting Bag \#2: $P(\#2)=\frac{1}{2}$
            \item The probability of selecting Bag a blue marble given that Bag \#2 was chosen:$P(B|\#2)=\frac{5}{10}$
            \item The probability of selecting a blue marble. We'd have to select Bag \#1 or Bag \#2 and then select a blue marble: $P(B)=\frac{1}{2}\cdot \frac{5}{10}+\frac{1}{2}\cdot 1=\frac{3}{4}$. When the case arises that you can't find this, use the other version of Bayes' Theorem.
        \end{itemize}
        Now we plug and play: $$P(\#2|B)=\frac{P(B|\#2)P(\#2)}{P(B)}=\frac{\frac{5}{10} \cdot (\frac{1}{2})}{\frac{3}{4}}=\frac{1}{3}$$
        Note that we could have used the other version of Bayes' Theorem:
        \[ P(\#2|B)=\frac{P(B|\#2)P(\#2)}{P(B|\#2)P(\#2)+P(B|\overline{\#2})P(\overline{\#2})}=\frac{\frac{5}{10} \cdot (\frac{1}{2})}{\frac{5}{10} \cdot (\frac{1}{2})+1 \cdot (\frac{1}{2})}=\frac{1}{3} \]
    \end{solution}

    
\question The zombie apocalypse has come. However, it's not as bad as the movies. Only \%1 of the population will get infected. 90\% of the people who are infected with this disease will test positive, but \%5 of the people tested will have false positives (meaning the test shows positive but they are \textbf{not} infected). What is the probability that a person is infected if the test shows positive? Approximate your answer to the $4^{\text{th}}$ decimal place.
    \begin{solution}
        We are looking for the probability that someone is infected given that they have tested positive, e.g., $P(I|T)$. As this is a conditional probability questions, we should immediately consider using Bayes' theorem:
         \[ P(I|T)=\displaystyle \frac{P(T|I)P(I)}{P(T|I)P(I)+P(T|\overline{I})P(\overline{I})} \]
        So we need the following:
        \begin{itemize}
            \item The probability that a person is infected: $P(I)=.01$
            \item The probability that an infected person tests positive (true positive): $P(T|I)=.90$
            \item The probability that a person is not infected: $P(\overline{I})=.99$
            \item The probability that a non-infected person tests positive (false positive): $P(T|\overline{I})=.05$ 
        \end{itemize}
        
        Now we plug these values into the formula:
         \[ P(I|T)=\displaystyle \frac{P(T|I)P(I)}{P(T|I)P(I)+P(T|\overline{I})P(\overline{I})}
         =\frac{(.90)(.01)}{(.90)(.01)+(.05)(.99)}\approx 0.1538
         \]
        Note that using this version of Bayes' theorem makes sense because we were given the probabilities of a false positive and a true positive. To use the other version,
        $$P(I|T)=\frac{P(T|I)P(I)}{P(T)}$$
        we would need to find $P(T)$, but this is the probability that someone who is infected tests positive OR someone who is not infected tests positive, i.e.,
        $$P(T)=P(T\cap I)+P(T\cap \overline{I})=P(T|I)P(I)+P(T|\overline{I})P(\overline{I})$$
        The denominator of the other version. Though more complicated to write, it also applies to more general situations.
    \end{solution}


\end{questions}

\end{document}
